\section{La cuenta inversa: dieciséis patrones}\label{sec:inversa}

Para obtener el valor del límite necesitamos evaluar el denominador de la doble
cuenta. Vista desde la capa con una reina más, una reparación se reconoce por el
par de reinas que dejó en el tablero. Por eso el contador inverso sigue pares y
su normalización es de orden $n^2$.

En esta sección $0\le h\le K$ (la capa cero está incluida), el proceso es $\QQ_h$ y $w$ es una
función continua y acotada en $[0,1]^2$, evaluada en los centros de las raíces. Se estudia
$I_w$ sobre la colocación final.

\subsection{Patrones y compensador explícito}

Escribimos un par ordenado compatible como $e_1=(r,j)$, $e_2=(i,c)$: se inserta $e^\star=(i,j)$
y se crea la raíz $o=(r,c)$. Sus ocho líneas son distintas, y las dos diagonales $d_1,d_2$ de
$e^\star$ son distintas de ellas; sólo se suman pares compatibles. Para
$L\subseteq\{1,2\}$ y $J\subseteq\{1,2\}$ definimos
\[
 Z_{L,J}=\prod_{g\in L}O_{e_g}\prod_{g\notin L}\chi_{e_g}\prod_{i\in J}\chi_{d_i}\prod_{i\notin J}V_{d_i}.
\]
Son dieciséis patrones por par; sea $\ell=|L|$. Para $z_t\ge\zeta$, con
$\alpha=(Nz_t)^{-1}$, $\beta_i=p_{d_i}/(Nz_t)$ y $k_g=\pi_{e_g}/(Nz_t^4)$, el generador ideal es
\[
 \cL Z_{L,J}=-\bigl[4(2-\ell)+|J|\bigr]\alpha\,Z_{L,J}-\sum_{i\notin J}\beta_i\,Z_{L,J\cup\{i\}}
 +\sum_{g\in L}k_g\,Z_{L\setminus\{g\},J}.
\]
Un nacimiento no borra las otras marcas, por la geometría anterior; la parte real de una muerte
de $V$ exige la marca activa y por eso aumenta $J$. En la región~\eqref{eq:regularidad}, con tolerancia
$\eps n$, el
resto respecto del generador verdadero satisface
\[
 |R_{L,J}|\le C_\zeta\Bigl[(\eps/n+n^{-2})Z_{L,J}+(\eps/n^2+n^{-3})\sum_{g\in L}Z_{L\setminus\{g\},J}\Bigr];
\]
la cota $O(n^{-2})$ de los borrados dobles incluye tanto las reales, de codegrado uno, como las
virtuales cartesianas.

Fijamos un horizonte determinista $T\le T_\zeta$ y ponemos $\Delta y=(T-t)/N$,
$\theta_g=\pi_{e_g}\Delta y/z_t^4$ y $\xi_i=p_{d_i}\Delta y/z_t$.

\begin{lema}[Compensador]\label{lem:compensador}
La función
\[
 \Phi^{(e_1,e_2)}_t=\bigl(O_{e_1}+\theta_1\chi_{e_1}\bigr)\bigl(O_{e_2}+\theta_2\chi_{e_2}\bigr)
 \bigl(V_{d_1}-\xi_1\chi_{d_1}\bigr)\bigl(V_{d_2}-\xi_2\chi_{d_2}\bigr)
\]
resuelve $\partial_t\Phi+\cL\Phi=0$ con valor terminal
$O_{e_1}O_{e_2}V_{d_1}V_{d_2}$ en $t=T$. Sus coeficientes
$c_{L,J}=\prod_{g\notin L}\theta_g\prod_{i\in J}(-\xi_i)$ en la base $Z_{L,J}$ satisfacen
$|c_{L,J}|\le C_\zeta n^{-(2-\ell)}$, $|\partial_tc_{L,J}|\le C_\zeta n^{-(3-\ell)}$ y
$|\partial_t^2c_{L,J}|\le C_\zeta n^{-(4-\ell)}$.
\end{lema}

\begin{proof}
El generador ideal es la suma de operadores que actúan sobre una sola de las diez marcas
distintas, de modo que basta que cada factor resuelva su propia ecuación hacia atrás. Para
$O_{e_g}+\theta_g\chi_{e_g}$: $\cL O_{e_g}=k_g\chi_{e_g}$ y $\cL\chi_{e_g}=-4\alpha\chi_{e_g}$, y
$\theta_g'+k_g-4\alpha\theta_g=0$ porque $\partial_t\Delta y=\partial_tz_t=-1/N$. Para
$V_{d_i}-\xi_i\chi_{d_i}$: $\cL V_{d_i}=-\beta_i\chi_{d_i}$ y $\cL\chi_{d_i}=-\alpha\chi_{d_i}$, y
$\xi_i'+\beta_i-\alpha\xi_i=0$. En $t=T$, $\Delta y=0$. Las cotas de los coeficientes siguen de
$\pi_{e_g}=O(n^{-1})$ y $p_{d_i}=O(1)$.
\end{proof}

Multiplicando cada par por $w(\bar r,\bar c)$ y sumando, el compensador agregado vale en $t=0$,
donde $\chi\equiv V\equiv1$ y $O\equiv0$,
\[
 \Phi_0=y_T^2\sum_{(e_1,e_2)}\pi_{e_1}\pi_{e_2}\bigl(1-y_Tp_{d_1}\bigr)\bigl(1-y_Tp_{d_2}\bigr)
 w(\bar r,\bar c),\qquad y_T=T/N .
\]
De nuevo es la solución de un sistema lineal, no una factorización de la ley real.

\subsection{Concentración}

\begin{lema}\label{lem:conc_inversa}
Para toda tolerancia fija $\gamma>0$,
$\Pr\{|I_w(T)-\Phi_0|>\gamma n^2,\ T<\sigma'\}\le e^{-c_{\zeta,\gamma,w}n}$, y la unión sobre
$O(n)$ horizontes conserva la cola.
\end{lema}

\begin{proof}
Para $\ell=0,1,2$, la suma de los indicadores con $\ell$ aristas ocupadas es $O(n^{4-\ell})$: hay
a lo sumo $n$ reinas por arista ocupada y $n^2$ casillas posibles por arista disponible. En un paso, la
suma de cambios absolutos es $O(n^{3-\ell})$: una diagonal fija corta cada fila y cada columna
una sola vez, lo que afecta a $O(n)$ pares con dos aristas ocupadas, $O(n^2)$ con una y $O(n^3)$
con ninguna; para los indicadores $\chi_{e_g}$ se usa que cada línea incide en $O(n)$ casillas,
y los nacimientos dan los mismos órdenes. La tabla siguiente reúne estas cuentas con los
coeficientes del Lema~\ref{lem:compensador} y los restos del generador; en un nacimiento se cuenta
el estado fuente, con $\ell-1$ aristas ocupadas, y no el destino.
\begin{center}\small
\begin{tabular}{@{}cccccc@{}}
\toprule
$\ell$ & patrones & coeficiente & cambio por paso & resto de muertes & resto de nacimientos\\
\midrule
$0$ & $O(n^4)$ & $O(n^{-2})$ & $O(n^3)$ & $n^4\cdot n^{-2}\,(\eps/n+n^{-2})$ & ---\\
$1$ & $O(n^3)$ & $O(n^{-1})$ & $O(n^2)$ & $n^3\cdot n^{-1}\,(\eps/n+n^{-2})$
 & $n^4\cdot n^{-1}\,(\eps/n^2+n^{-3})$\\
$2$ & $O(n^2)$ & $O(1)$ & $O(n)$ & $n^2\cdot(\eps/n+n^{-2})$ & $n^3\cdot(\eps/n^2+n^{-3})$\\
\bottomrule
\end{tabular}
\end{center}
En cada fila, el coeficiente por el cambio por paso es $O(n)$, lo que acota los saltos, y cada
resto es $O(\eps n+1)$ por paso. Los segundos
incrementos temporales y los términos mixtos suman $O(1)$ por paso, por las cotas de derivadas
del Lema~\ref{lem:compensador} y la cuenta de cambios. Se congelan estado y tiempo en la parada
$\sigma'$ (Lema~\ref{lem:conc_directa}), incluida la transición de salida, cuya probabilidad
hasta $T_\zeta$ es $e^{-cn}$. La deriva acumulada es $O_\zeta(\eps n^2+n)$ y los saltos de la
martingala son $O_\zeta(n)$ durante $O(n)$ pasos; Azuma, eligiendo primero la tolerancia y
después $\eps$, da la cota. No hace falta unir eventos sobre las $n^4$ parejas.
\end{proof}
